% Einheit 1 – Wurzelziehen und Lösbarkeit – Hauptnummern 12 bis 23
\einheitenkopf[e1][1 Wurzelziehen]{Einheit 1 von 4 · Wurzelziehen und Lösbarkeit}
\zweigzeile{Hier lernst du, Gleichungen wie $x^2 = c$ durch Wurzelziehen zu lösen und zu sagen, ob es zwei, eine oder keine Lösung gibt · neu in diesem Jahr, je nach Buch bis Klasse 10 (am Gymnasium schon in Klasse 8 oder 9) · P10 · baut auf: Quadrieren, Wurzel ziehen, lineare Gleichungen, Scheitelpunktform}
\setcounter{aufgabe}{11}

\begin{aufgabe}{Ich erkenne, ob eine Gleichung quadratisch ist.}
Kreuze an: quadratisch oder linear? Löse nichts.
\begin{teile}
\teil $x^2 = 50$ \hfill \kreuz{quadratisch} \kreuz{linear}
\teil $5x - 3 = 12$ \hfill \kreuz{quadratisch} \kreuz{linear}
\teil $(x + 2)^2 = 1$ \hfill \kreuz{quadratisch} \kreuz{linear}
\teil $x \cdot x = 16$ \hfill \kreuz{quadratisch} \kreuz{linear}
\teil $3x + x = 12$ \hfill \kreuz{quadratisch} \kreuz{linear}
\end{teile}
\end{aufgabe}

\begin{aufgabe}{Ich erkenne, ob rechts eine positive Zahl, null oder eine negative Zahl steht.}
Sieh nur auf die rechte Seite und kreuze an. Rechne nichts.
\begin{teile}
\teil $x^2 = 64$ \hfill \kreuz{positiv} \kreuz{null} \kreuz{negativ}
\teil $x^2 = -4$ \hfill \kreuz{positiv} \kreuz{null} \kreuz{negativ}
\teil $2x^2 = 0$ \hfill \kreuz{positiv} \kreuz{null} \kreuz{negativ}
\teil $x^2 = 2{,}5$ \hfill \kreuz{positiv} \kreuz{null} \kreuz{negativ}
\teil $x^2 = -0{,}01$ \hfill \kreuz{positiv} \kreuz{null} \kreuz{negativ}
\end{teile}
\end{aufgabe}

\verfahren{Gleichungen der Form $x^2 = c$}
\begin{aufgabe}{Ich kann eine Gleichung $x^2 = c$ durch Wurzelziehen lösen.}
Löse die Gleichung. Gib beide Lösungen an.
\rechenplatz{2}
\begin{gleichungsraster}[2]
\gl{x^2 = 16} & \gl{x^2 = 81} \\
\gl{x^2 = 4} & \gl{x^2 = 144} \\
\anweisung{Lass die Wurzel stehen und gib zusätzlich Näherungswerte auf zwei Stellen nach dem Komma an.}
\gl{x^2 = 10} & \\
\anweisung{Hat die Gleichung keine Lösung, schreibe das als Satz.}
\gl{x^2 = -16} & \\
\end{gleichungsraster}
\end{aufgabe}

\begin{aufgabe}{Ich kann $x^2$ erst allein stellen und dann die Wurzel ziehen.}
Löse die Gleichung.
\rechenplatz{3}
\begin{gleichungsraster}[3]
\gl{x^2 - 4 = 12} & \gl{x^2 + 3 = 67} \\
\gl{2x^2 = 50} & \gl{3x^2 + 1 = 13} \\
\gl{2x^2 + 9 = 1} & \gl{4x^2 + 3 = 3} \\
\end{gleichungsraster}
\end{aufgabe}

\verfahren{Gleichungen mit quadrierter Klammer}
\begin{aufgabe}{Ich kann eine Gleichung mit quadrierter Klammer rückwärts lösen.}
Löse die Gleichung.
\rechenplatz{3}
\begin{gleichungsraster}[3]
\gl{(x - 2)^2 = 9} & \gl{(x - 1)^2 = 25} \\
\gl{(x - 5)^2 = 4} & \gl{(x - 3)^2 = 1} \\
\gl{(x + 2)^2 = 25} & \gl{(x + 5)^2 = 0} \\
\end{gleichungsraster}
\end{aufgabe}

\verfahren{Lösungen prüfen und zählen}
\begin{aufgabe}{Ich kann prüfen, ob eine Zahl eine quadratische Gleichung löst.}
Setze die Zahl ein und kreuze an: wahre Aussage (wA) oder falsche Aussage (fA).
\begin{teile}
\teil $x = 7$ \quad in \quad $x^2 + 1 = 50$ \hfill \kreuz{wA} \kreuz{fA}
\teil $x = -7$ \quad in \quad $x^2 + 1 = 50$ \hfill \kreuz{wA} \kreuz{fA}
\teil $x = 3$ \quad in \quad $(x - 1)^2 = 9$ \hfill \kreuz{wA} \kreuz{fA}
\teil $x = -2$ \quad in \quad $(x + 5)^2 = 9$ \hfill \kreuz{wA} \kreuz{fA}
\teil $x = -1$ \quad in \quad $3x^2 - 5 = -8$ \hfill \kreuz{wA} \kreuz{fA}
\end{teile}
\end{aufgabe}

\begin{aufgabe}{Ich kann am Graphen ablesen, wie viele Lösungen eine Gleichung hat.}
Die Parabel gehört zu $y = (x + 1)^2$. Lies für jede Gleichung ab, wie viele Lösungen sie hat, und gib die Lösungen an.

\begin{ksys}[xmin=-5,xmax=4,ymin=-3,ymax=6,ablesen]
\parabel{1}{-1}{0}{}
\gerade[3]{0}{4}{y=4}
\gerade[3]{0}{1}{y=1}
\gerade[3]{0}{-2}{y=-2}
\end{ksys}

\begin{teile}
\teil $(x + 1)^2 = 4$ \hfill Anzahl: \leerfeld \quad Lösungen: \leerfeld
\teil $(x + 1)^2 = 1$ \hfill Anzahl: \leerfeld \quad Lösungen: \leerfeld
\teil $(x + 1)^2 = 0$ \hfill Anzahl: \leerfeld \quad Lösungen: \leerfeld
\teil $(x + 1)^2 = -2$ \hfill Anzahl: \leerfeld \quad Lösungen: \leerfeld
\end{teile}
\end{aufgabe}

\begin{aufgabe}{Ich kann Gleichungen mit keiner, einer oder zwei Lösungen angeben.}
Ergänze die rechte Seite so, dass die Gleichung die angegebene Zahl von Lösungen hat.
\begin{teile}
\teil genau eine Lösung: \quad $(x - 3)^2 = {}$ \leerfeld
\teil zwei Lösungen: \quad $(x + 2)^2 = {}$ \leerfeld
\teil keine Lösung: \quad $x^2 + 5 = {}$ \leerfeld
\teil keine Lösung: \quad $3x^2 = {}$ \leerfeld
\anweisung{Schreibe selbst eine Gleichung.}
\teil eine Gleichung mit quadrierter Klammer, die keine Lösung hat: \leerfeld
\end{teile}
\end{aufgabe}

\begin{aufgabe}{Ich kann Aussagen über die Lösungen einer Gleichung beurteilen.}
Kreuze an, ob die Aussage wahr oder falsch ist.
\begin{teile}
\teil \sachtabelle{lcc}{Aussage & wahr & falsch}{„$(x - 6)^2 = 0$ hat genau eine Lösung.“ & $\square$ & $\square$ \\ „$2$ und $10$ sind Lösungen von $(x + 6)^2 = 16$.“ & $\square$ & $\square$}
\teil Gib eine quadratische Gleichung an, die keine Lösung hat: \leerfeld \hfill (P10 2021 OS)
\end{teile}
\end{aufgabe}

\begin{aufgabe}{Ich finde den Fehler beim Wurzelziehen.}
Finde den Fehler, benenne ihn und korrigiere die Rechnung.
\begin{teile}
\teil Paul rechnet: \rechnung{x^2 &= 64 \\ x &= 8}
\teil Lea rechnet: \rechnung{(x + 9)^2 &= 4 \\ x + 9 &= 2 \quad \text{oder} \quad x + 9 = -2 \\ x_1 &= 11 \quad \text{oder} \quad x_2 = 7}
\teil Ben rechnet: \rechnung{x^2 &= -81 \\ x &= -9}
\end{teile}
\end{aufgabe}

\begin{aufgabe}{Ich kann begründen, wie viele Lösungen eine Gleichung hat.}
Begründe ohne zu rechnen.
\begin{teile}
\teil Warum hat $x^2 = 20$ zwei Lösungen?
\teil Warum hat $x^2 = -20$ keine Lösung?
\teil Warum hat $(x - 9)^2 = 0$ nur eine Lösung?
\end{teile}
\end{aufgabe}

\begin{aufgabe}{Ich kann mit Wurzelziehen eine Länge berechnen.}
Berechne. Runde auf zwei Stellen nach dem Komma.
\begin{teile}
\teil Ein quadratisches Beet hat eine Fläche von $121\,\text{m}^2$. Wie lang ist eine Seite? \\ Seitenlänge: \leerfeld[m]
\teil Eine quadratische Tischplatte hat eine Fläche von $1{,}44\,\text{m}^2$. Wie lang ist eine Seite? \\ Seitenlänge: \leerfeld[m]
\teil Eine zylinderförmige Dose soll $10\,\text{cm}$ hoch sein und $500\,\text{cm}^3$ fassen. Für das Volumen gilt $V = \pi \cdot r^2 \cdot h$. Wie groß muss der Radius $r$ sein? (P10 2022 OS) \\ $r = {}$ \leerfeld[cm]
\end{teile}

\zylinder{1.2}{2}{r}{h}

\end{aufgabe}
